\documentclass[11pt,spanish]{article}

% =========================================================
% PLANTILLA BASE PARA INFORMES ACADÉMICOS / TÉCNICOS
% =========================================================

% --- Idioma y codificación ---
\usepackage[spanish,es-tabla]{babel}
\usepackage[T1]{fontenc}
\usepackage[utf8]{inputenc}

% --- Márgenes / papel ---
\usepackage[
    left=2.5cm,
    right=2.5cm,
    top=2.5cm,
    bottom=2.5cm,
    headsep=0.5cm,
    footskip=0.8cm
]{geometry}

% --- Matemáticas ---
\usepackage{amsmath}

% --- Tablas, enlaces y elementos visuales ---
\usepackage{array}
\usepackage{tabularx}
\usepackage{booktabs}
\usepackage{longtable}
\usepackage{graphicx}
\usepackage{float}
\usepackage{xcolor}
\usepackage{tikz}
\usetikzlibrary{positioning,shapes.geometric,arrows.meta,calc}
\usepackage{enumitem}
\usepackage{ragged2e}
\usepackage[hidelinks]{hyperref}

% --- Encabezado y pie de página ---
\usepackage{fancyhdr}
\usepackage{lastpage}

% --- Interlineado ---
\linespread{1.2}
\sloppy

% =========================================================
% DATOS DEL DOCUMENTO
% =========================================================
\newcommand{\Institucion}{Instituto Tecnológico de Costa Rica}
\newcommand{\Campus}{Campus Tecnológico Central Cartago}
\newcommand{\DocumentoTitulo}{Plan de Comunicaciones}
\newcommand{\DocumentoSubtitulo}{Plataforma Universitaria Integrada: Modernización y Centralización de Procesos Académicos y Administrativos}
\newcommand{\Curso}{Administración de proyectos}
\newcommand{\Profesor}{Alicia Marcela Salazar Hernández}
\newcommand{\FechaDocumento}{9 de junio de 2026}
\newcommand{\PeriodoAcademico}{I Semestre}
\newcommand{\AnioDocumento}{2026}

\newcommand{\EstudiantesPortada}{%
    Mauricio González Prendas: 2024143009\\
    Isaac Jiménez Blanco: 2024202804\\
    Tamara Robles Camacho: 2024099342%
}

% Datos que aparecen en el encabezado.
\newcommand{\DocHeaderLeft}{}
\newcommand{\DocHeaderRight}{Plan de Comunicaciones}
\newcommand{\DocFooterRight}{Página~\thepage\ de~\pageref{LastPage}}

% Metadatos PDF.
\title{\DocumentoTitulo}
\author{\EstudiantesPortada}
\date{\FechaDocumento}

% =========================================================
% CONFIGURACIÓN DE TABLAS Y UTILIDADES
% =========================================================
\newcolumntype{Y}{>{\RaggedRight\arraybackslash}X}
\renewcommand{\arraystretch}{1.22}
\setlength{\LTpre}{0.5em}
\setlength{\LTpost}{0.5em}
\setlist[itemize]{leftmargin=*, label=\textbullet}
\setcounter{tocdepth}{3}


% Imagen segura: si el archivo no existe, muestra un recuadro de marcador.
\newcommand{\SafeImage}[2][0.92\textwidth]{%
    \IfFileExists{#2}{%
        \includegraphics[width=#1]{#2}%
    }{%
        \fbox{%
            \begin{minipage}[c][4cm][c]{#1}
            \centering
            Imagen pendiente:\\
            \texttt{#2}
            \end{minipage}%
        }%
    }%
}

% Figura reutilizable.
\newcommand{\EvidenceFigure}[3]{%
    \begin{figure}[H]
    \centering
    \SafeImage{#1}
    \caption{#2}
    \label{#3}
    \end{figure}%
}

% =========================================================
% ENCABEZADO Y PIE
% =========================================================
\pagestyle{fancy}
\fancyhf{}
\fancyhead[L]{\DocHeaderLeft}
\fancyhead[R]{\DocHeaderRight}
\renewcommand{\headrulewidth}{0.4pt}
\renewcommand{\footrulewidth}{0.4pt}
\fancyfoot[R]{\DocFooterRight}

% Estilo equivalente para páginas horizontales.
\fancypagestyle{widefancy}{%
    \fancyhf{}%
    \fancyhead[L]{\DocHeaderLeft}%
    \fancyhead[R]{\DocHeaderRight}%
    \renewcommand{\headrulewidth}{0.4pt}%
    \renewcommand{\footrulewidth}{0.4pt}%
    \fancyfoot[R]{\DocFooterRight}%
}

% =========================================================
% PÁGINAS HORIZONTALES PARA TABLAS ANCHAS
% =========================================================
\makeatletter
\newcommand{\SyncPageBuilderDimensions}{%
    \global\vsize=\textheight
    \global\@colroom=\textheight
    \global\@colht=\textheight
}
\makeatother

\newenvironment{widepage}{%
    \clearpage
    \hypersetup{pageanchor=false}%
    \paperwidth=297mm
    \paperheight=210mm
    \pdfpagewidth=297mm
    \pdfpageheight=210mm
    \newgeometry{left=2.5cm,right=2.5cm,top=2.5cm,bottom=2.5cm,headsep=0.5cm,footskip=0.8cm}%
    \setlength{\textwidth}{243mm}%
    \setlength{\textheight}{156mm}%
    \setlength{\linewidth}{\textwidth}%
    \setlength{\columnwidth}{\textwidth}%
    \hsize=\textwidth
    \setlength{\headwidth}{\textwidth}%
    \SyncPageBuilderDimensions
    \pagestyle{widefancy}%
}{%
    \clearpage
    \paperwidth=210mm
    \paperheight=297mm
    \pdfpagewidth=210mm
    \pdfpageheight=297mm
    \restoregeometry
    \SyncPageBuilderDimensions
    \hypersetup{pageanchor=true}%
    \pagestyle{fancy}%
}

% =========================================================
% PORTADA
% =========================================================
\newcommand{\MakeCover}{%
\begin{titlepage}
\thispagestyle{empty}
\centering

{\bfseries \huge \Institucion \par}
\vspace{0.25cm}
{\LARGE \Campus \par}

\vfill

{\huge \DocumentoTitulo \par}


\vfill

{\bfseries \LARGE Profesora: \par}
{\Large \Profesor \par}

\vfill

{\bfseries \LARGE Estudiantes: \par}
{\Large \EstudiantesPortada \par}

\vfill

{\bfseries\LARGE Fecha: \par}
{\Large \FechaDocumento \par}

\vfill

{\LARGE \PeriodoAcademico \par}
{\LARGE \AnioDocumento \par}

\end{titlepage}%
}

% =========================================================
% DOCUMENTO
% =========================================================
\begin{document}
\hypersetup{pageanchor=false}

% =========================
% Portada
% =========================
\MakeCover

% =========================
% Índice
% =========================
\pagenumbering{roman}
\pagestyle{empty}
\addtocontents{toc}{\protect\thispagestyle{empty}}
\tableofcontents
\clearpage
\pagenumbering{arabic}
\hypersetup{pageanchor=true}
\pagestyle{fancy}

% =========================
% Contenido
% =========================
\section{Información general del documento}

\begin{table}[H]
\centering
\begin{tabularx}{\textwidth}{p{5cm} Y}
\toprule
\textbf{Nombre del proyecto} & Plataforma Universitaria Integrada \\
\textbf{Organización} & Universidad Tecnológica La Mejor \\
\textbf{Documento} & \DocumentoTitulo \\
\textbf{Project manager} & Tamara Robles Camacho \\
\textbf{Fecha} & \FechaDocumento \\
\bottomrule
\end{tabularx}
\end{table}

\section{Introducción}

El presente documento corresponde al Entregable 12, Plan de Comunicaciones, del proyecto Plataforma Universitaria Integrada, desarrollado para la Universidad Tecnológica La Mejor. Este plan forma parte de la planificación del proyecto y tiene como finalidad definir cómo se comunicará la información entre los integrantes del equipo, el Sponsor y los stakeholders relacionados con el proyecto.

La comunicación es un elemento clave dentro de la administración de proyectos, ya que permite mantener alineadas a las partes interesadas, reducir confusiones, coordinar tareas, informar avances, documentar decisiones y atender riesgos o cambios de manera ordenada. En este proyecto, la comunicación adquiere mayor importancia porque el equipo es pequeño, los roles son combinados y el plazo de entrega es corto. La entrega final debe realizarse antes del 24 de junio de 2026, por lo que cualquier atraso, cambio o falta de coordinación puede afectar directamente el alcance, el cronograma, la calidad o los entregables finales.

El proyecto consiste en diseñar y desarrollar un prototipo Front-End web navegable de una Plataforma Universitaria Integrada. Esta plataforma representa procesos académicos, administrativos y financieros mediante módulos como Login, Dashboard principal, Portal Estudiantil, Portal Docente, Portal Administrativo, Portal Financiero y Dashboard Ejecutivo. Además del prototipo, el proyecto incluye entregables de gestión como Business Case, Project Charter, Registro de Stakeholders, Definición del Alcance, Cronograma del Proyecto, Organigrama, Plan de Calidad, Plan de Comunicaciones, Registro de Riesgos, Control de Cambios y documentos de cierre.

Este Plan de Comunicaciones se relaciona directamente con los entregables elaborados previamente. El Project Charter define los roles principales del equipo y la necesidad de mantener una coordinación constante. El Registro de Stakeholders identifica a los interesados internos y externos. La Definición del Alcance establece las restricciones del proyecto, incluyendo que las decisiones relevantes deben comunicarse al equipo para evitar documentos incongruentes. El Cronograma del Proyecto define las actividades, fechas, responsables e hitos. El Organigrama muestra la estructura de comunicación entre Sponsor, Project Manager, equipo ejecutor y stakeholders externos.

Para este proyecto, GitHub será utilizado como repositorio central del prototipo y como espacio compartido para organizar entregables, evidencias y versiones finales. WhatsApp se utilizará como canal rápido para coordinación diaria entre integrantes. Zoom se utilizará para reuniones de seguimiento o revisiones cuando sea necesario. El correo electrónico se utilizará para comunicaciones formales, principalmente con el Sponsor o stakeholders externos cuando se requiera dejar evidencia clara de una decisión, validación o entrega.

El objetivo principal de este plan es asegurar que cada persona reciba la información correcta, en el momento adecuado y por el canal correspondiente. De esta forma, se busca evitar pérdida de información, duplicidad de trabajo, confusiones sobre responsabilidades, cambios no documentados y problemas de trazabilidad.

\section{Propósito del plan de comunicaciones}

El propósito del Plan de Comunicaciones es establecer las reglas, canales, responsables, frecuencia y formatos mediante los cuales se compartirá la información del proyecto Plataforma Universitaria Integrada.

Este documento permite responder las siguientes preguntas:

\begin{itemize}
    \item Qué información debe comunicarse.
    \item A quién debe comunicarse.
    \item Quién es responsable de comunicarla.
    \item Con qué frecuencia debe comunicarse.
    \item Qué canal debe utilizarse.
    \item Qué formato debe tener la comunicación.
    \item Qué información debe quedar registrada como evidencia.
    \item Cómo se comunicarán riesgos, cambios, problemas y avances.
\end{itemize}

El plan busca que la comunicación del proyecto sea clara, oportuna, trazable y adecuada para cada stakeholder. No toda comunicación tendrá el mismo nivel de formalidad. Algunas coordinaciones internas podrán realizarse por WhatsApp, mientras que decisiones importantes, cambios de alcance, entregas finales o validaciones deberán quedar documentadas en GitHub, correo electrónico o minuta de reunión.

\section{Objetivos del plan de comunicaciones}

Los objetivos del Plan de Comunicaciones son los siguientes:

\subsection{Objetivo 1}

Garantizar que el equipo de trabajo mantenga una comunicación constante sobre avances, tareas, dudas, riesgos, cambios y entregables pendientes.

\subsection{Objetivo 2}

Definir canales oficiales para evitar que la información importante se pierda en conversaciones informales o quede sin respaldo.

\subsection{Objetivo 3}

Mantener informados al Sponsor y stakeholders externos sobre los avances relevantes del proyecto, especialmente cuando se requiera validación o revisión de entregables.

\subsection{Objetivo 4}

Asegurar que los documentos y entregables del proyecto se mantengan organizados en un espacio compartido, utilizando GitHub como repositorio central del proyecto.

\subsection{Objetivo 5}

Reducir el riesgo de inconsistencias entre documentos, especialmente en temas como alcance, cronograma, costos, calidad, riesgos, roles y fechas.

\subsection{Objetivo 6}

Establecer un mecanismo formal para comunicar cambios que puedan afectar alcance, tiempo, costo, calidad o riesgos del proyecto.

\section{Stakeholders considerados}

Este plan considera los stakeholders internos y externos definidos en los entregables anteriores.

\subsection{Stakeholders internos}

\begin{table}[H]
\centering
\caption{Stakeholders internos}
\label{tab:stakeholders-internos}
{\small
\begin{tabularx}{\textwidth}{p{3.7cm} p{4.1cm} Y}
\toprule
\textbf{Nombre} & \textbf{Rol} & \textbf{Necesidad de comunicación} \\
\midrule
Tamara Robles Camacho & Project Manager & Coordinar avances, revisar entregables, comunicar decisiones y dar seguimiento al cronograma \\
Isaac Jiménez Blanco & Analista de Negocio y Documentador 1 & Mantener trazabilidad documental, comunicar riesgos, cambios, calidad y entregables \\
Jorge Abarca Quiros & Analista de Negocio y Documentador 2 & Elaborar y revisar entregables documentales, apoyar la gestión de riesgos, costos y comunicaciones \\
Mauricio González Prendas & Desarrollador Front-End 1 y QA & Comunicar avance técnico de portales básicos, problemas del prototipo, pruebas y correcciones \\
Carlos Sainz Arce & Desarrollador Front-End 2 y QA & Comunicar avance técnico de portales financiero y ejecutivo, y pruebas de calidad \\
\bottomrule
\end{tabularx}
}
\end{table}

\subsection{Stakeholders externos}

\begin{table}[H]
\centering
\caption{Stakeholders externos}
\label{tab:stakeholders-externos}
{\small
\begin{tabularx}{\textwidth}{p{3.7cm} p{3.9cm} Y}
\toprule
\textbf{Nombre} & \textbf{Rol} & \textbf{Necesidad de comunicación} \\
\midrule
Dr. Roberto Mora Fallas & Sponsor & Recibir avances generales, validar entregables principales y aprobar decisiones relevantes \\
Lic. Patricia Vargas Soto & Directora Administrativa & Validar necesidades administrativas representadas en el prototipo \\
Ing. Carlos Brenes Mora & Director de TI & Revisar aspectos técnicos generales y viabilidad visual del prototipo \\
Prof. Andrea Solís Zamora & Representante Docente & Validar claridad del Portal Docente \\
Estudiantes & Usuarios finales & Apoyar validaciones básicas del Portal Estudiantil si se requiere \\
Departamento Financiero & Área financiera & Validar representación del Portal Financiero \\
Proveedor de Hosting & Proveedor externo & Brindar referencia técnica si se requiere publicar o alojar evidencia del prototipo \\
\bottomrule
\end{tabularx}
}
\end{table}

\section{Canales oficiales de comunicación}

Para mantener orden en la comunicación del proyecto se definen los siguientes canales:

\subsection{GitHub}

Será el canal principal para centralizar el repositorio del prototipo Front-End, documentos finales, evidencias y versiones controladas. También puede utilizarse para registrar issues, tareas técnicas, defectos y cambios relacionados con el desarrollo.

\subsection{WhatsApp}

Será utilizado para comunicación rápida entre los integrantes del equipo. Este canal servirá para coordinar horarios, resolver dudas rápidas, avisar avances inmediatos y solicitar revisión de tareas. No será considerado canal formal para aprobar cambios de alcance, presupuesto o cronograma.

\subsection{Zoom}

Será utilizado para reuniones de seguimiento, revisión de entregables, validación de avances y resolución de dudas importantes. Cuando una reunión produzca decisiones relevantes, estas deberán documentarse mediante una minuta breve o mensaje formal posterior.

\subsection{Correo electrónico}

Será utilizado para comunicaciones formales con el Sponsor o stakeholders externos. También se utilizará cuando se requiera dejar evidencia de una aprobación, entrega, solicitud de revisión o decisión relevante.

\subsection{Microsoft Project}

Será utilizado para comunicar y respaldar la planificación del cronograma, incluyendo tareas, fechas, responsables, costos, hitos, dependencias y ruta crítica. Las actualizaciones relevantes del cronograma deberán comunicarse al equipo.

\subsection{Microsoft Word y PDF}

Serán utilizados para preparar y entregar documentos formales del proyecto. Los PDF se usarán como evidencia final de los entregables.

\section{Reglas generales de comunicación}

Para evitar confusión y mantener trazabilidad, el equipo aplicará las siguientes reglas:

\begin{itemize}
    \item Toda decisión que afecte alcance, tiempo, costo, calidad o riesgos deberá documentarse formalmente.
    \item WhatsApp podrá usarse para coordinación rápida, pero no para aprobar cambios importantes.
    \item Los entregables finales deberán guardarse en GitHub o en una carpeta compartida asociada al repositorio.
    \item Cada documento deberá tener un nombre claro y relacionado con el entregable correspondiente.
    \item Las reuniones importantes deberán dejar una minuta breve o resumen de acuerdos.
    \item Los cambios al prototipo deberán registrarse en GitHub mediante commits o issues.
    \item Los problemas técnicos o defectos relevantes deberán comunicarse al equipo y documentarse.
    \item La Project Manager será responsable de dar seguimiento general a la comunicación del proyecto.
    \item El Analista de Negocio y Documentador deberá revisar que los documentos mantengan coherencia entre sí.
    \item El Desarrollador Front-End y QA deberá comunicar avances técnicos, defectos y correcciones.
    \item Los stakeholders externos serán consultados únicamente cuando se requiera validación general de procesos o módulos.
\end{itemize}

\section{Matriz de comunicaciones}

La siguiente matriz define qué información se comunicará, quién será responsable, a quién se dirigirá, con qué frecuencia, por qué canal y con qué nivel de prioridad.

\begin{widepage}
{\scriptsize
\setlength{\tabcolsep}{2.4pt}
\begin{longtable}{p{3.05cm} p{3.15cm} p{3.65cm} p{2.6cm} p{2.55cm} p{3.0cm} p{1.65cm}}
\caption{Matriz de comunicaciones}\label{tab:matriz-comunicaciones}\\[-0.5em]
\toprule
\textbf{Información a comunicar} & \textbf{Responsable} & \textbf{Destinatarios} & \textbf{Frecuencia} & \textbf{Canal} & \textbf{Formato} & \textbf{Prioridad} \\
\midrule
\endfirsthead
\toprule
\textbf{Información a comunicar} & \textbf{Responsable} & \textbf{Destinatarios} & \textbf{Frecuencia} & \textbf{Canal} & \textbf{Formato} & \textbf{Prioridad} \\
\midrule
\endhead
Avance general del proyecto & Project Manager & Equipo y Sponsor & Semanal o según avance & Zoom o correo & Resumen breve de avance & Alta \\
Estado de tareas del equipo & Project Manager & Equipo interno & Cada 2 o 3 días & WhatsApp & Mensaje de seguimiento & Media \\
Entregables documentales finalizados & Analista de Negocio y Documentador & Equipo interno & Por evento & GitHub y WhatsApp & Documento en Word o PDF & Alta \\
Versiones finales de documentos & Analista de Negocio y Documentador & Equipo interno y Sponsor si aplica & Por evento & GitHub o correo & Archivo final & Alta \\
Avance del prototipo Front-End & Desarrollador Front-End y QA & Equipo interno & Cada 2 o 3 días & WhatsApp y GitHub & Mensaje y commit & Alta \\
Cambios en el código del prototipo & Desarrollador Front-End y QA & Equipo interno & Por evento & GitHub & Commit o issue & Media \\
Defectos encontrados en el prototipo & Desarrollador Front-End y QA & Equipo interno & Por evento & GitHub o WhatsApp & Issue o mensaje & Alta \\
Riesgos identificados & Analista de Negocio y Documentador & Project Manager y equipo & Por evento & WhatsApp y documento formal & Registro de riesgos & Alta \\
Cambios solicitados & Project Manager o responsable del cambio & Equipo y Sponsor si aplica & Por evento & Correo, Zoom o documento formal & Solicitud de cambio & Alta \\
Evaluación de cambios & Analista de Negocio y Documentador & Project Manager y equipo & Por evento & Documento formal y reunión & Evaluación de impacto & Alta \\
Validación de módulos & Desarrollador Front-End y QA & Equipo interno y stakeholders si aplica & Según avance & Zoom o WhatsApp & Revisión visual o checklist & Media \\
Validación del Portal Docente & Project Manager & Prof. Andrea Solís Zamora y equipo & Por evento & Correo o Zoom & Solicitud de revisión & Media \\
Validación del Portal Financiero & Project Manager & Departamento Financiero y equipo & Por evento & Correo o Zoom & Solicitud de revisión & Media \\
Validación administrativa & Project Manager & Lic. Patricia Vargas Soto y equipo & Por evento & Correo o Zoom & Solicitud de revisión & Media \\
Revisión técnica general & Project Manager & Ing. Carlos Brenes Mora y equipo & Por evento & Correo o Zoom & Solicitud de revisión & Media \\
Entrega final del proyecto & Project Manager & Sponsor, equipo y profesora & Al cierre & GitHub y correo si aplica & Carpeta final y documentos PDF & Alta \\
\bottomrule
\end{longtable}
}
\end{widepage}

\section{Comunicación según prioridad}

Para manejar adecuadamente la información, se establecen tres niveles de prioridad.

\subsection{Prioridad alta}

Corresponde a información que puede afectar alcance, cronograma, costo, calidad, riesgos o entrega final. Debe comunicarse de inmediato al equipo y documentarse formalmente.

\noindent Ejemplos:

\begin{itemize}
    \item Cambio en el alcance del prototipo.
    \item Atraso en una tarea crítica.
    \item Error crítico en el Front-End.
    \item Inconsistencia importante entre documentos.
    \item Modificación del presupuesto.
    \item Riesgo que pueda afectar la entrega final.
\end{itemize}

\noindent Canales recomendados:

\begin{itemize}
    \item WhatsApp para alerta inicial.
    \item Zoom para análisis si es necesario.
    \item Documento formal, GitHub o correo para dejar evidencia.
\end{itemize}

\subsection{Prioridad media}

Corresponde a información necesaria para coordinar trabajo, revisar avances o resolver dudas, pero que no afecta de forma inmediata el cumplimiento del proyecto.

\noindent Ejemplos:

\begin{itemize}
    \item Revisión de un documento.
    \item Solicitud de apoyo en redacción.
    \item Corrección visual del prototipo.
    \item Ajuste menor en una tabla.
    \item Coordinación de reunión.
\end{itemize}

\noindent Canales recomendados:

\begin{itemize}
    \item WhatsApp.
    \item GitHub si se relaciona con código o archivos.
    \item Zoom si requiere discusión.
\end{itemize}

\subsection{Prioridad baja}

Corresponde a información operativa o administrativa que no afecta directamente el alcance ni la fecha de entrega.

\noindent Ejemplos:

\begin{itemize}
    \item Confirmación de horario.
    \item Comentarios menores de formato.
    \item Organización de archivos.
    \item Recordatorios simples.
\end{itemize}

\noindent Canales recomendados:

\begin{itemize}
    \item WhatsApp.
    \item Mensaje breve al equipo.
\end{itemize}

\section{Gestión de reuniones}

Las reuniones del proyecto se realizarán únicamente cuando sean necesarias para revisar avances, resolver dudas o validar entregables importantes. Debido al tiempo limitado del proyecto, se dará prioridad a reuniones cortas y enfocadas.

\subsection{Tipos de reuniones}

\subsubsection{Reunión de seguimiento interno}

\noindent Objetivo: revisar avance de tareas, problemas y próximos pasos.\\
Participantes: equipo interno.\\
Frecuencia: una o dos veces por semana, según necesidad.\\
Canal: Zoom o llamada grupal.\\
Responsable: Project Manager.

\subsubsection{Reunión de revisión de entregables}

\noindent Objetivo: revisar documentos importantes antes de considerarlos finalizados.\\
Participantes: equipo interno.\\
Frecuencia: por evento.\\
Canal: Zoom.\\
Responsable: responsable del entregable.

\subsubsection{Reunión de validación con stakeholders}

\noindent Objetivo: solicitar retroalimentación general sobre un módulo o documento.\\
Participantes: Project Manager, responsable del módulo y stakeholder correspondiente.\\
Frecuencia: según necesidad.\\
Canal: Zoom o correo.\\
Responsable: Project Manager.

\subsubsection{Reunión de cierre}

\noindent Objetivo: revisar que los entregables estén completos y listos para entrega final.\\
Participantes: equipo interno.\\
Frecuencia: una vez al final del proyecto.\\
Canal: Zoom.\\
Responsable: Project Manager.

\subsection{Minuta de reunión}

Cuando una reunión genere acuerdos importantes, se deberá dejar un resumen breve con:

\begin{itemize}
    \item Fecha de reunión.
    \item Participantes.
    \item Temas tratados.
    \item Acuerdos principales.
    \item Responsables asignados.
    \item Fecha de seguimiento.
\end{itemize}

\section{Gestión de documentos y entregables}

La gestión de documentos es parte fundamental del Plan de Comunicaciones, ya que el proyecto incluye varios entregables administrativos y técnicos.

Los documentos y evidencias del proyecto se organizarán en GitHub, utilizando una estructura clara de carpetas.

\subsection{Reglas para documentos}

\begin{itemize}
    \item Cada archivo debe tener un nombre claro.
    \item Los documentos finales deberán guardarse en PDF.
    \item Las versiones editables podrán mantenerse en Word cuando sea necesario.
    \item El archivo de Microsoft Project deberá conservarse en formato .mpp.
    \item Las evidencias importantes deberán guardarse en PDF o imagen.
    \item No se deben subir versiones duplicadas sin identificación clara.
    \item Los documentos finales deberán revisarse antes de ser compartidos.
\end{itemize}

\subsection{Ejemplos de nombres de archivo}

\begin{itemize}
    \item Entregable 2 - Project Charter.pdf
    \item Entregable 4 - Definición del Alcance.pdf
    \item Entregable 6 - Cronograma del Proyecto.mpp
    \item Entregable 11 - Plan de Calidadpdf
    \item Entregable 12 - Plan de Comunicaciones.pdf
\end{itemize}

\section{Comunicación de riesgos, cambios y problemas}

Los riesgos, cambios y problemas deben comunicarse de forma clara porque pueden afectar directamente el cumplimiento del proyecto.

\subsection{Comunicación de riesgos}

Cuando se identifique un riesgo, se deberá comunicar al equipo y registrar en el documento correspondiente.

La comunicación debe incluir:

\begin{itemize}
    \item Descripción del riesgo.
    \item Causa posible.
    \item Impacto esperado.
    \item Probabilidad estimada.
    \item Responsable.
    \item Acción de respuesta.
\end{itemize}

Ejemplo: Si se detecta que el desarrollo del Front-End puede atrasarse, Mauricio González Prendas deberá comunicarlo al equipo. La Project Manager revisará el impacto en cronograma y el Analista de Negocio y Documentador actualizará el registro de riesgos si corresponde.

\subsection{Comunicación de cambios}

Cualquier cambio que afecte alcance, tiempo, costo, calidad o riesgos debe ser tratado formalmente. No se aprobarán cambios importantes únicamente por WhatsApp.

La comunicación de un cambio debe incluir:

\begin{itemize}
    \item Nombre del solicitante.
    \item Descripción del cambio.
    \item Justificación.
    \item Impacto en alcance.
    \item Impacto en cronograma.
    \item Impacto en costos.
    \item Impacto en calidad.
    \item Recomendación.
    \item Decisión final.
\end{itemize}

\subsection{Comunicación de problemas}

Un problema es una situación que ya ocurrió y que puede afectar el proyecto. Debe comunicarse lo antes posible.

\noindent Ejemplos:

\begin{itemize}
    \item No se puede abrir el archivo de Microsoft Project.
    \item Hay error crítico en la navegación del prototipo.
    \item Un documento contradice el alcance.
    \item No se logra subir evidencia al repositorio.
    \item El equipo no tiene claridad sobre una tarea.
\end{itemize}

Los problemas de prioridad alta deberán notificarse por WhatsApp de inmediato y luego documentarse en GitHub, correo o el entregable correspondiente.

\section{Responsables de comunicación}

La comunicación del proyecto tendrá responsables definidos según el tipo de información.

\subsection{Project Manager}

Tamara Robles Camacho

\begin{itemize}
    \item Coordinar la comunicación general del proyecto.
    \item Organizar reuniones de seguimiento.
    \item Comunicar avances generales al Sponsor.
    \item Validar que el equipo esté informado sobre decisiones relevantes.
\end{itemize}

\subsection{Analista de Negocio y Documentador 1}

Isaac Jiménez Blanco

\begin{itemize}
    \item Mantener coherencia entre documentos.
    \item Comunicar avances de entregables documentales (mitad).
    \item Registrar riesgos, cambios y evaluaciones.
    \item Verificar que la información documentada sea trazable.
\end{itemize}

\subsection{Analista de Negocio y Documentador 2}

Jorge Abarca Quiros

\begin{itemize}
    \item Comunicar avances de entregables documentales (mitad).
    \item Apoyar la gestión y comunicación de riesgos, costos y comunicaciones.
    \item Colaborar en el mantenimiento de la documentación del proyecto.
\end{itemize}

\subsection{Desarrollador Front-End 1 y QA}

Mauricio González Prendas

\begin{itemize}
    \item Comunicar avances técnicos del prototipo (módulos básicos).
    \item Registrar cambios relevantes en GitHub.
    \item Reportar defectos encontrados.
    \item Informar correcciones realizadas.
    \item Comunicar problemas técnicos que puedan afectar el cronograma.
\end{itemize}

\subsection{Desarrollador Front-End 2 y QA}

Carlos Sainz Arce

\begin{itemize}
    \item Comunicar avances técnicos del prototipo (módulos financiero y ejecutivo).
    \item Reportar defectos de calidad y correcciones en su sección.
    \item Colaborar en las pruebas de calidad técnica del front-end.
\end{itemize}

\subsection{Sponsor}

Dr. Roberto Mora Fallas

\begin{itemize}
    \item Recibir información de avance general.
    \item Revisar entregables principales.
    \item Aprobar decisiones relevantes cuando afecten el propósito del proyecto.
    \item Validar el cierre del proyecto.
\end{itemize}

\subsection{Stakeholders externos}

\begin{itemize}
    \item Brindar retroalimentación general cuando se les consulte.
    \item Validar módulos relacionados con su área.
    \item Comunicar observaciones que mejoran la representación del prototipo.
\end{itemize}

\section{Seguimiento y control de las comunicaciones}

El seguimiento de las comunicaciones permite verificar que la información fluya correctamente durante el proyecto. Para ello, se aplicarán las siguientes acciones:

\begin{itemize}
    \item Revisar periódicamente que los entregables estén en la ubicación compartida definida.
    \item Confirmar que las decisiones importantes queden documentadas.
    \item Verificar que los integrantes del equipo conozcan sus tareas y fechas.
    \item Revisar que los cambios importantes pasen por el proceso formal correspondiente.
    \item Confirmar que los riesgos relevantes sean comunicados y registrados.
    \item Validar que las versiones finales estén correctamente nombradas.
    \item Revisar que los acuerdos de reunión sean entendidos por los participantes.
\end{itemize}

\subsection{Indicadores simples de seguimiento}

\begin{table}[H]
\centering
\caption{Indicadores simples de seguimiento}
\label{tab:indicadores-seguimiento}
\begin{tabularx}{\textwidth}{Y p{3cm}}
\toprule
\textbf{Indicador} & \textbf{Meta} \\
\midrule
Entregables importantes almacenados en GitHub & 100\% \\
Cambios relevantes documentados formalmente & 100\% \\
Riesgos importantes comunicados al equipo & 100\% \\
Reuniones con acuerdos documentados cuando aplique & 100\% \\
Documentos finales con nombre correcto & 100\% \\
Integrantes informados sobre tareas críticas & 100\% \\
\bottomrule
\end{tabularx}
\end{table}

Si durante el seguimiento se detecta que un canal no está funcionando, la Project Manager deberá proponer un ajuste. Por ejemplo, si WhatsApp genera pérdida de información, las decisiones deberán trasladarse a correo, GitHub o minuta de reunión.

\section{Conclusión}

El Plan de Comunicaciones del proyecto Plataforma Universitaria Integrada establece la forma en que se compartirá la información entre el equipo, el Sponsor y los stakeholders externos. Este documento permite definir canales, responsables, frecuencia, formatos y reglas para que la comunicación sea clara, ordenada y trazable.

La comunicación será especialmente importante debido a que el proyecto tiene un plazo corto, varios entregables documentales, un prototipo Front-End navegable y roles combinados dentro del equipo. Por esta razón, se utilizarán canales distintos según el tipo de información. GitHub funcionará como repositorio central del proyecto, WhatsApp como canal rápido de coordinación, Zoom como medio para reuniones y correo electrónico como canal formal para validaciones o decisiones importantes.

Este plan también ayuda a reducir riesgos de mala comunicación, como documentos incongruentes, cambios no registrados, pérdida de evidencia, duplicidad de trabajo o atrasos en tareas críticas. Al definir reglas claras, el equipo puede mantener mejor control sobre el alcance, el cronograma, los costos, la calidad, los riesgos y los cambios.

Con la aplicación de este Plan de Comunicaciones, el proyecto tendrá una estructura más ordenada para compartir información, documentar decisiones y asegurar que los entregables se mantengan alineados con lo definido en el Business Case, Project Charter, Registro de Stakeholders, Definición del Alcance, Cronograma, Organigrama y Plan de Calidad.

\end{document}
